\section*{Capítulo \ref{ch:g6:exploring-our-environment} --- Explorar nuestro entorno}

\begin{solution}{exo:g6:exploring-our-environment:1}
\omterm{def:g6:exploring-our-environment:components}{Seres vivos} (musgo, cochinillas); \omterm{def:g6:exploring-our-environment:components}{materia mineral} (la piedra del muro,
el charco); huellas y restos de \omterm{def:g6:exploring-our-environment:components}{seres vivos} (\omterm{def:g1:plants-around-us:parts}{hojas} muertas, la \omterm{def:g1:animals-around-us:covering}{concha}
de caracol).
\end{solution}

\begin{solution}{exo:g6:exploring-our-environment:2}
La \omterm{def:g6:exploring-our-environment:conditions}{temperatura}, con un termómetro (en grados Celsius); la \omterm{def:g6:exploring-our-environment:conditions}{humedad} del
aire y de la tierra; la \omterm{def:g6:exploring-our-environment:conditions}{luz}, desde el pleno sol a la sombra profunda.
\end{solution}

\begin{solution}{exo:g6:exploring-our-environment:3}
\omterm{def:g6:exploring-our-environment:components}{Seres vivos}: el helecho, el ciempiés. Mineral: el charco, la piedra
del muro. Huellas: la \omterm{def:g1:animals-around-us:covering}{concha} de caracol, los excrementos de pájaro.
\end{solution}

\begin{solution}{exo:g6:exploring-our-environment:4}
Dieciséis grados (\qty{38}{\degreeCelsius} frente a
\qty{22}{\degreeCelsius}) y pleno sol frente a sombra profunda (y
también seco frente a húmedo). La hierba de sequía y resplandor encaja
con el pie sur; el musgo, los \omterm{prop:g4:plants-without-seeds:damp}{helechos} y las cochinillas encajan con
el rincón fresco y húmedo.
\end{solution}

\begin{solution}{exo:g6:exploring-our-environment:5}
Su cubierta deja escapar el agua, así que al sol y en seco su cuerpo
pierde el agua mortalmente deprisa. Por eso las cochinillas se
encuentran en \omterm{def:g6:exploring-our-environment:microhabitat}{microhábitats} húmedos y oscuros --- bajo piedras, bajo
madera muerta.
\end{solution}

\begin{solution}{exo:g6:exploring-our-environment:6}
Por ejemplo: una avellana roída --- el ratón de campo; una telaraña
--- su araña; unas pisadas en el barro --- la garza (los excrementos,
las \omterm{def:g1:animals-around-us:covering}{plumas} y las \omterm{def:g1:animals-around-us:covering}{conchas} sirven igual).
\end{solution}

\begin{solution}{exo:g6:exploring-our-environment:7}
Una caja cerrada, con una mitad de papel húmedo y otra de papel seco,
y las cochinillas soltadas a lo largo de la línea del centro; unos
minutos después casi todas están en el lado húmedo (y lo mismo con la
sombra frente a la \omterm{def:g6:exploring-our-environment:conditions}{luz}). Demuestra que cada \omterm{def:g1:animals-around-us:animal}{animal} se instala donde
las \omterm{def:g6:exploring-our-environment:conditions}{condiciones} le van a su cuerpo --- el mapa del grupo se dibuja
solo, sin ningún rebaño.
\end{solution}

\begin{solution}{exo:g6:exploring-our-environment:8}
Porque solo se pueden comparar medidas hechas de la misma manera: una
diferencia tiene que venir de los sitios y no de la medición. Es la
versión de campo de la prueba justa --- varía una cosa (el sitio) y
todo lo demás se mantiene igual.
\end{solution}

\begin{solution}{exo:g6:exploring-our-environment:9}
Encima de la piedra: iluminado, seco, cálido --- líquenes (con la
mosca que toma el sol de visita). Debajo de la piedra: oscuro, húmedo,
fresco --- cochinillas (o el ciempiés). A cinco centímetros, dos
conjuntos de habitantes bien emparejados.
\end{solution}

\begin{solution}{exo:g6:exploring-our-environment:10}
Paso 1: marcar un metro cuadrado concreto. Paso 2: medir --- la
lectura del termómetro anotada, la \omterm{def:g6:exploring-our-environment:conditions}{humedad} de la tierra notada y
apuntada, la clase de \omterm{def:g6:exploring-our-environment:conditions}{luz} anotada. Paso 3: rebuscar como es debido,
levantando las piedras y volviéndolas a poner, y enumerar las tres
clases de componentes. Paso 4: anotar allí mismo --- recuentos, no
``algunos''. Paso 5: con un solo sitio no había comparación posible;
reconocer un segundo sitio de la misma manera.
\end{solution}

\begin{solution}{exo:g6:exploring-our-environment:11}
Explicación: la hiedra prospera en el intervalo de \omterm{def:g6:exploring-our-environment:conditions}{condiciones} del
muro norte (sombra, \omterm{def:g6:exploring-our-environment:conditions}{humedad}) y no en el del muro sur (resplandor,
sequía). Prueba: medir el pie de los dos muros --- \omterm{def:g6:exploring-our-environment:conditions}{temperatura},
\omterm{def:g6:exploring-our-environment:conditions}{humedad}, \omterm{def:g6:exploring-our-environment:conditions}{luz} --- y comprobarlo con las preferencias conocidas de la
hiedra; o buscar hiedra en otros muros y anotar qué orientación
comparten esos muros.
\end{solution}

\begin{solution}{exo:g6:exploring-our-environment:12}
Quitar la diferencia de la pared: hacer la elección en una caja
redonda (con la pared igual en todas partes) y con las mitades húmeda
y seca intercambiadas entre pasadas --- húmeda al este en la pasada 1,
húmeda al oeste en la pasada 2. Si las cochinillas siguen la \omterm{def:g6:exploring-our-environment:conditions}{humedad}
allí donde se ponga, las paredes quedan descartadas; entre las pasadas
solo cambió la condición que se estaba probando, y todo lo demás se
mantuvo igual.
\end{solution}

\begin{solution}{pb:g6:exploring-our-environment:1}
\textbf{1.} El seto alto: da sombra a la cuneta N (y la abriga),
mientras que la cuneta S mira al sol todo el día.

\textbf{2.} \omterm{def:g6:exploring-our-environment:conditions}{Temperatura}: \qty{19}{\degreeCelsius} frente a
\qty{31}{\degreeCelsius}. \omterm{def:g6:exploring-our-environment:conditions}{Humedad}: la tierra de la cuneta N está
húmeda (se pega) y la de la cuneta S está seca (polvo). La \omterm{def:g6:exploring-our-environment:conditions}{luz}
completa el trío: sombra frente a pleno sol.

\textbf{3.} Las \omterm{def:g6:exploring-our-environment:conditions}{condiciones} cambian a lo largo del día: a las 6 de la
tarde cualquier sitio está más fresco y con menos \omterm{def:g6:exploring-our-environment:conditions}{luz}, así que la
diferencia medida mezclaría el sitio con la hora. La misma hora
mantiene justa la comparación.

\textbf{4.} Para que el esfuerzo de búsqueda sea igual: más superficie
o más tiempo en un sitio encontraría allí más habitantes --- una
prueba injusta de qué cuneta es más rica.

\textbf{5.} Las dos listas nombran solo \omterm{def:g6:exploring-our-environment:components}{seres vivos}. Faltan: los
componentes minerales (tierra, piedras, agua) y las huellas --- que un
inventario en regla también tiene que anotar.

\textbf{6.} Cochinillas: cuneta N --- necesitan \omterm{def:g6:exploring-our-environment:conditions}{humedad} y sombra o se
secan. El sapo: cuneta N --- \omterm{def:g1:the-five-senses:sense}{piel} húmeda, necesita un refugio húmedo.
La lagartija: cuneta S --- calienta su cuerpo al sol y caza allí donde
puede tomarlo.

\textbf{7.} O bien viven caracoles allí escondidos (quizá activos solo
en las horas húmedas de la noche), o bien la \omterm{def:g1:animals-around-us:covering}{concha} es un resto de
otro sitio o de otra estación --- una huella demuestra presencia
\emph{en algún momento}, no residencia ahora.

\textbf{8.} Oscuro, húmedo, fresco. La regla de la
\cref{prop:g6:exploring-our-environment:decide}: cada punto está
habitado por las \omterm{def:g5:biodiversity:species}{especies} cuyos intervalos encajan con sus \omterm{def:g6:exploring-our-environment:conditions}{condiciones}.

\textbf{9.} N tiene musgo y sapos porque su sombra la mantiene fresca
y húmeda, que es lo que sus cuerpos necesitan; S tiene lagartijas y
hierbas de \omterm{def:g1:plants-around-us:parts}{hojas} cerosas porque su pleno sol la hace caliente y seca,
que es lo que le va a una cazadora que toma el sol y a una \omterm{def:g1:plants-around-us:plant}{planta} a
prueba de sequía.

\textbf{10.} Sin el seto, la cuneta N recibe pleno sol: más caliente,
más seca --- unas \omterm{def:g6:exploring-our-environment:conditions}{condiciones} que van derivando hacia las de la cuneta
S de hoy. Es de esperar que los amantes de la \omterm{def:g6:exploring-our-environment:conditions}{humedad} (musgo,
\omterm{prop:g4:plants-without-seeds:damp}{helechos}, cochinillas, el sapo) menguen y que se instalen los amantes
del sol (hierbas secas, saltamontes, quizá la lagartija).

\textbf{11.} Porque el nuevo reconocimiento tiene que diferenciarse
del primero en una sola cosa --- la ausencia del seto. Unos sitios,
instrumentos o estación nuevos añadirían diferencias propias, y el
cambio ya no se le podría achacar al seto: la regla de la prueba justa
a escala de paisaje.

\textbf{12.} Cada cochinilla da vueltas y gira hasta que su cuerpo
está a gusto, así que el grupo acaba junto allí donde está la \omterm{def:g6:exploring-our-environment:conditions}{humedad}
--- las posiciones de los \omterm{def:g1:animals-around-us:animal}{animales} registran las \omterm{def:g6:exploring-our-environment:conditions}{condiciones} con tanta
seguridad como las lecturas de un termómetro. Nadie las lleva en
rebaño; el mapa lo dibuja el bienestar de cada individuo, que es
exactamente por lo que los habitantes se pueden leer como medidas.
\end{solution}
